\section{Related Work}

\paragraph{Single Object Generation.}
Building on the success of text-to-image generation models~\cite{rombach2022high,ramesh2022hierarchical,kang2023scaling,yu2022scaling}, there has been rapid progress in 3D generative models conditioned on text or images. A prominent class of methods leverages 2D diffusion priors for 3D generation~\cite{poole2022dreamfusion,wang2023score,lin2023magic3d,chen2023fantasia3d,metzer2023latent,wang2024prolificdreamer,sun2023dreamcraft3d,long2023wonder3d,michel2022text2mesh}, with DreamFusion~\cite{poole2022dreamfusion} introducing Score Distillation Sampling (SDS) to optimize 3D representations using gradients from 2D diffusion models.
Subsequent works have extended SDS with multi-view diffusion models, improving both 3D generation quality and single-view reconstruction~\cite{liu2023zero123,wang2023imagedream,shi2023mvdream,liu2023one2345,zhou2023sparsefusion,liu2023one2345++,long2023wonder3d,shi2023zero123++,liu2023syncdreamer}. Another research direction trains large-scale transformers to generate 3D shapes in a feed-forward manner~\cite{hong2023lrm,li2023instant3d,xu2023dmv3d,TripoSR2024}, relying on curated, large-scale 3D asset datasets.
While these models can generate visually compelling shapes, they often ignore the physical properties, such as stability, of the object, which are essential for real-world applications. To address this, recent efforts have incorporated physics-based simulation into the generation pipeline to produce self-supporting 3D objects by optimizing physical attributes such as mass distribution~\cite{guo2024physically,chen2024atlas3d,yan2024phycage,cai2024gaussian}. Additionally, PhysDreamer~\cite{zhang2024physdreamer}, optimizes physical properties like Young's modulus and initial velocity to generate dynamic motions that are both visually plausible and physically grounded, guided by video diffusion priors.


\paragraph{Scene Generation.} 
While single-object generation methods produce visually appealing assets, they often lack scale awareness and spatial grounding, making scene composition challenging.
The primary bottlenecks in 3D scene generation include decomposing scenes into individual assets, estimating their relative scale and pose, and ensuring physical feasibility (e.g., contact, stability).
Several works address these challenges through multi-stage pipelines.
Early methods such as~\cite{vilesov2023cg3d,chen2024comboverse,han2024reparo} adopt object-centric reconstruction followed by layout and geometry optimization using physical constraints or differentiable rendering.
\citet{shriram2024realmdreamer} lifts the scene image to 3D point clouds as a whole, inpaints occluded regions, and refines appearance using 2D diffusion priors.
Recently, Large Language Models (LLMs) and VLMs are increasingly leveraged to infer spatial relations and scene structure.
\citet{gao2024graphdreamer} constructs a scene graph with objects as nodes and their relations as edges.
\citet{zhou2024gala3d} uses a VLM to generate a coarse layout, which is subsequently refined with rendering losses and physical constraints.
\citet{yao2025cast} infers a scene graph describing simplified pairwise relationships between objects, and use them to optimize object poses and scales.
\citet{wang2025embodiedgengenerative3dworld} similarly leverages LLM-based reasoning to query objects’ relative sizes and physical properties, enabling more plausible scene layouts. However, none of these methods can ensure physically accurate contact handling or maintain physical stability in the generated scene.
Scene-level optimization under SDS loss is commonly used for joint geometry-text alignment~\cite{zhou2025layoutdreamer}, while \citet{huang2024midi} and \citet{xu2024comp4d} explore multi-instance and 4D compositional generation guided by spatial and trajectory priors.
Other approaches incorporate physical property estimation into 3D representations, either from visual cues~\cite{zhao2024automated} or explicit user input~\cite{chen2025physgen3d}, to support dynamic simulation or interaction~\cite{liu2024physgen}.
Recent systems such as Blender-MCP and \citet{li2025phip} integrate LLM reasoning and generative priors into graphics tools like Blender, enabling fine-grained control and interactive refinement. \citet{sun2024layoutvlm} propose LayoutVLM, which uses a VLMs to generate differentiable spatial relations and jointly optimize 3D layouts in indoor scenes. However, most existing scene generation methods focus primarily on layout composition. They either omit physical reasoning altogether or incorporate only simple physics priors to prevent object interpenetration, without modeling accurate contact interactions or ensuring physically stable scene layouts.
We thus address this gap by novelly augmenting text-to-3D scene generation with differentiable rigid body contact simulation.

